\begin{center}\bfseries Résumé\end{center}
\phantomsection
\addcontentsline{toc}{section}{Résumé}
\begingroup
\fontsize{9.5}{10.7}\selectfont
\setlength{\parskip}{1.5pt plus .5pt minus .5pt}
% Composición local: conserva resumen y palabras clave en la apertura.
\fontsize{10.2}{11.6}\selectfont
\setlength{\parskip}{1pt plus .5pt minus .5pt}
\enlargethispage{2\baselineskip}
\noindent
On construit une réponse constitutive positive et inversible du vide
électromagnétique : deux canaux orientés déterminent conjointement les
coordonnées adimensionnelles de permittivité, de perméabilité, d’impédance et de
vitesse de propagation. On démontre leurs relations exactes et la
récupération unique de la paire angulaire qui les produit. La même réponse
permet de reconstruire la forme de l’ellipse d’action et d’exprimer la
fonctionnelle de Barbero--Immirzi au moyen de ses deux valeurs propres étiquetées.
Le résultat relie ainsi la détermination des valeurs à
l’information structurelle conservée par leur évaluation.

La construction part d’un noyau formel mathématique appelé
holographie modulaire triadique, ci-après HMT. \emph{All Possible Paths}
(APP) est la structure arithmétique des chemins sur le support toroïdal
$(\mathbb Z/9\mathbb Z)^2$ ; ses déplacements cardinaux forment le
graphe de Cayley additif, et ses feuillets conservent les évaluations somme et
produit avec résidu et quotient. TRIT est la signature orientée
$\{-1,0,+1\}$ qui distingue les régimes locaux. Le
\emph{Topological Phase Kernel} (TPK) compose la sélection, le transport,
la mise à jour de mémoire et le retour sur cet état enrichi. Le
prolongement nonadique et la clôture dodécaphasique engendrent les coordonnées
corrélées et les registres dont proviennent l’action, la paire
angulaire et les incidences utilisées par la réponse.

Les secteurs $90$ et $120$ du transport $T$ fixent
$\mathcal V=(I-T^{90})(I-T^{120})^{-1}$. Sa fonction scalaire
$f(s)=(1+s+s^2)/(1+s+s^2+s^3)$ est une bijection décroissante de $(0,1)$
sur $(3/4,1)$ : cette propriété établit l’inversion cubique et son
domaine. L’échange de feuillets permute la permittivité et la perméabilité,
inverse l’impédance et conserve la vitesse. Dans la même coordinatisation dimensionnelle,
on démontre en outre l’identité entre la vitesse constitutive et la
vitesse cinématique de l’horloge.

La réalisation sur le cycle à douze phases apporte le défaut de trace
$-1/12$, une représentation de $f$ comme quotient de déterminants et un quart de tour
de carré $-I$. La double intégration logarithmique de sa résolvante
produit le moment de Catalan. Dans la section de Pell $2-\sqrt3$, on
démontre une identité exacte avec la constante de Catalan et
$\log(2+\sqrt3)$ ; les deux arguments constitutifs maintiennent leur
détermination angulaire propre. L’inversion de $f$ compose ensuite la
réponse du vide avec la fonctionnelle orientée de Barbero--Immirzi.
La restitution intégrale démontre également que la fonction constitutive
détermine le moment orienté complet sous la condition
$\mathcal C_2(s)\to0$ à l’origine ; sa limite en $s=1$ retrouve
la constante de Catalan. On conserve ainsi une correspondance fonctionnelle
avec son orientation et ses arguments, outre ses évaluations.

La polarisation discrète s’obtient à partir du bilan borné des vacances,
et le préfacteur électronique $\sqrt3/4$ à partir de la norme du commutateur des
projections associées aux fibres de somme et de produit du support
arithmétique tridimensionnel. La section
positive de charge, construite avec $\alpha$, l’action et l’impédance,
détermine les relations de von Klitzing, de conductance, de flux magnétique
et de Josephson, y compris leurs identités croisées et leur covariance sous
le retour de l’action. Les réalisations dimensionnelles conservent les
sections d’action, de charge et de vitesse et leurs changements d’unités. L’évaluation
numérique correspond aux coordonnées internes ;
sa comparaison physique conserve explicitement la représentation et la
coordinatisation dimensionnelle utilisées.

La séparation positive du mode commun et du mode orienté spécifie
l’information conservée par la normalisation unimodulaire. La résolution
de Fourier de l’incidence dodécaphasique identifie quatre modes et
leurs poids exacts, en distinguant le défaut de rang, les multiplicités
signées et le coefficient singulier du retour.

L’inversion induit également une loi associative de composition des
réponses, avec une coordonnée angulaire additive. Dans la réalisation
électrofaible à l’arbre, les lecteurs de jauge et de Higgs retrouvent
conjointement les canaux et l’impédance, avec un domaine explicite.

\smallskip
\noindent\textbf{Mots-clés :} détermination structurelle des constantes ;
vide électromagnétique ; réponse constitutive ; inversion spectrale ;
holographie modulaire triadique ; orientation des feuillets ; constante de Catalan ;
unité de Pell ; Barbero--Immirzi ; polarisation ; quantification électrique.
\par\endgroup
